\bta{2011}



\begin{enumerate}
	%\renewcommand{\labelenumi}{\arabic{enumi}.}
	% A(\Alph) a(\alph) I(\Roman) i(\roman) 1(\arabic)
	%设定全局标号series=example	%引用全局变量resume=example
	%[topsep=-0.3em,parsep=-0.3em,itemsep=-0.3em,partopsep=-0.3em]
	%可使用leftmargin调整列表环境左边的空白长度 [leftmargin=0em]
	\item
	选择题 (本题满分 40 分，每小题 5 分)
	\begin{enumerate}
		%\renewcommand{\labelenumi}{\arabic{enumi}.}
		% A(\Alph) a(\alph) I(\Roman) i(\roman) 1(\arabic)
		%设定全局标号series=example	%引用全局变量resume=example
		%[topsep=-0.3em,parsep=-0.3em,itemsep=-0.3em,partopsep=-0.3em]
		%可使用leftmargin调整列表环境左边的空白长度 [leftmargin=0em]
		\item
		极限 $\lim\limits _{x \rightarrow + \infty}x\left[\left(1+\frac{1}{x}\right)^{x}-e\right]=$ \xzanswer{D} 
		
		
	\fourchoices
	{$ 0 $}
	{$\infty$}
	{$\frac{e}{2}$}
	{$-\frac{e}{2}$}

%	$ (1+\frac{1}{x})^{x}-e=e(e^{x\ln(1+\frac{1}{x} ) -1} -1) $，此时再用等价替换就好


	
	\item 
	若函数 $\left\{\begin{array}{ll}x^{2}+2 x \quad(x \geqslant 0) \\ \ln (1+a x) & (x<0)\end{array}\right.$ 在 $x=0$ 处可导, 则 $a$ 等于	 \xzanswer{B} 
	
	
\fourchoices
{$-2$}
{$ 2 $}
{$-1$}
{$ 1 $}

\item 
设极限 $\lim\limits _{x \rightarrow a} \frac{f(x)-f(a)}{(x-a)^{4}}=-2$, 则函数 $f(x)$ 在 $x=a$ 处 \xzanswer{C} 


\fourchoices
{导数值不等于$ 0 $}
{导数值等于$ 0 $ 但 $x=a$ 不是极值点}
{取得极大值}
{取得极小值}

\item 
设函数 $f(x)$ 在定义域内可导, $f(x)$ 的图像如下图所示, 则 $f^{\prime}(x)$ 的图像最有可能是 \xzanswer{D} 
\begin{figure}[h!]
	\centering
	\includesvg[width=0.23\linewidth]{picture/svg/601-001}
\end{figure}
\pfourchoices
{\includesvg[width=3.3cm]{picture/svg/601-002}}
{\includesvg[width=3.3cm]{picture/svg/601-003}}
{\includesvg[width=3.3cm]{picture/svg/601-004}}
{\includesvg[width=3.3cm]{picture/svg/601-005}}

\item 
设 $x^{2} \ln x$ 是函数 $f(x)$ 的一个原函数, 则不定积分 $\int x f^{\prime}(x) d x=$ \xzanswer{C} 


\fourchoices
{$\frac{2}{3} x^{3} \ln x+\frac{1}{9} x^{3}+C$}
{$2 x-x^{2} \ln x+C$}
{$x^{2} \ln x+x^{2}+C$}
{$3 x^{2} \ln x+x^{2}+C$}

\item 
设 $y_{1}, y_{2}$ 是二阶线性函数齐次微分方程 $y^{\prime \prime}+p(x) y^{\prime}+q(x) y=0$ 的两个特解, $C_{1}, C_{2}$ 是两个任意常数, 则 \xzanswer{B} 


\fourchoices
{$C_{1} y_{1}+C_{2} y_{2}$ 不一定是该微分方程的解}
{$C_{1} y_{1}+C_{2} y_{2}$ 是该微分方程的解, 但不一定是通解}
{$C_{1} y_{1}+C_{2} y_{2}$ 是该微分方程的解，但不是通解}
{$C_{1} y_{1}+C_{2} y_{2}$ 不是该微分方程的解}

\item 
若级数 $\sum\limits_{n=1}^{\infty} u_{n}$ 和 $\sum\limits_{n=1}^{\infty} v_{n}$ 都发散，则下列级数一定发散的是 \xzanswer{C} 


\fourchoices
{$\sum\limits_{n=1}^{\infty}\left(u_{n}+v_{n}\right)$}
{$\sum\limits_{n=1}^{\infty} u_{n} v_{n}$}
{$\sum\limits_{n=1}^{\infty}\left(\left|u_{n}\right|+\left|v_{n}\right|\right)$}
{$\sum\limits_{n=1}^{\infty}\left(u_{n}^{2}+v_{n}^{2}\right)$}

\item 
设两条抛物线 $y=n x^{2}+\frac{1}{n}$ 和 $y=(n+1) x^{2}+\frac{1}{n+1}$ 所围成的面积为 $A_{n}$, 则 $\lim\limits _{n \rightarrow \infty} A_{n}=$ \xzanswer{A}



\fourchoices
{$ 0 $}
{$ 1 $}
{$ 2 $}
{$ 3 $}


		
	\end{enumerate}


\banswer{
\begin{enumerate}
	%\renewcommand{\labelenumi}{\arabic{enumi}.}
	% A(\Alph) a(\alph) I(\Roman) i(\roman) 1(\arabic)
	\item
	\textbf{D}。由 $\ln(1+t)=t-\frac{t^2}{2}+o\left(t^2\right)$（取 $t=\frac1x\to0$）得
	\[
	x\ln\left(1+\frac1x\right)=1-\frac{1}{2x}+o\left(\frac1x\right),
	\]
	故 $\left(1+\frac1x\right)^{x}=e\cdot e^{-\frac{1}{2x}+o(1/x)}=e\left[1-\frac{1}{2x}+o\left(\frac1x\right)\right]$，
	于是
	\[
	x\left[\left(1+\frac1x\right)^{x}-e\right]\longrightarrow-\frac{e}{2}.
	\]

	\item
	\textbf{B}。两侧当 $x\to0$ 时均趋于 $0$，与 $f(0)=0$ 一致，故连续性自动成立。
	右导数 $f'_{+}(0)=\lim\limits_{x\to0^{+}}\frac{x^2+2x}{x}=2$；
	左导数 $f'_{-}(0)=\lim\limits_{x\to0^{-}}\frac{\ln(1+ax)}{x}=a$。
	可导要求两者相等，故 $a=2$。

	\item
	\textbf{C}。分母 $\to0$ 而比值有极限，必有 $f(x)-f(a)\to0$。
	又当 $x$ 充分接近 $a$（$x\neq a$）时 $\frac{f(x)-f(a)}{(x-a)^4}<-1<0$，
	而 $(x-a)^4>0$，故 $f(x)<f(a)$，即 $f$ 在 $x=a$ 处取极大值。
	（同时 $\frac{f(x)-f(a)}{x-a}=(x-a)^3\cdot\frac{f(x)-f(a)}{(x-a)^4}\to0$，
	故导数若存在必为 $0$，A、B 皆不对。）

	\item
	\textbf{D}。左支上 $f$ 单调上升且愈来愈陡，故 $f'>0$ 且 $f'$ 递增，
	$f'$ 的左支应整条位于 $x$ 轴上方——A、C 的左支穿过 $x$ 轴，排除。
	右支上 $f$ 自间断点右侧起急速上升，$f'(0^{+})$ 应是较大的正值，
	而 B 的右支从 $x$ 轴上出发，排除。
	再看 $f$ 在右支先取极大、后取极小，故 $f'$ 须由正变负（再变正），
	只有 D 的右支从正值出发并下穿 $x$ 轴。故选 D。

	\item
	\textbf{C}。由题设 $f(x)=\left(x^2\ln x\right)'=2x\ln x+x$，且 $\int f(x)dx=x^2\ln x+C$。
	分部积分
	\begin{align*}
	\int xf'(x)dx&=xf(x)-\int f(x)dx\\
	&=x\left(2x\ln x+x\right)-x^2\ln x+C\\
	&=x^2\ln x+x^2+C .
	\end{align*}

	\item
	\textbf{B}。由齐次线性方程解的叠加原理，$C_1y_1+C_2y_2$ 一定是解，故 A、D 错。
	但当 $y_1,y_2$ 线性相关时（例如 $y_2=2y_1$），$C_1y_1+C_2y_2$ 实质只含一个独立常数，
	不是通解；而当 $y_1,y_2$ 线性无关时它又确是通解。故只能说“是解，但不一定是通解”。

	\item
	\textbf{C}。若 $\sum|u_n|$ 收敛，则 $\sum u_n$ 绝对收敛从而收敛，与题设矛盾，
	故 $\sum|u_n|$ 发散；同理 $\sum|v_n|$ 发散。
	两个正项发散级数逐项相加，部分和仍 $\to+\infty$，故 C 一定发散。

	其余可用反例排除：取 $u_n=1,\ v_n=-1$，则 $\sum\left(u_n+v_n\right)=0$ 收敛，A 错；
	取 $u_n=v_n=\frac1n$（两者皆发散），则 $\sum u_nv_n=\sum\frac{1}{n^2}$、
	$\sum\left(u_n^2+v_n^2\right)=2\sum\frac1{n^2}$ 均收敛，B、D 错。

	\item
	\textbf{A}。令两式相等得 $x^2=\frac{1}{n(n+1)}$，记 $c_n=\frac{1}{\sqrt{n(n+1)}}$，
	则在 $\left[-c_n,c_n\right]$ 上两抛物线之差恰为 $c_n^2-x^2\geqslant0$，于是
	\[
	A_n=\int_{-c_n}^{c_n}\left(c_n^2-x^2\right)dx=\frac43c_n^3
	=\frac{4}{3\left[n(n+1)\right]^{3/2}} ,
	\]
	故 $\lim\limits_{n\to\infty}A_n=0$。
\end{enumerate}
}


\item
(本题满分 10 分)
设函数 $f(x)$ 在 $[0,2 \pi]$ 上有连续导数, 且 $f^{\prime}(x) \geqslant 0$, 求证：对任何正整数 $n$, 都有
\[
\left|\int_{0}^{2 \pi} f(x) \sin (nx) d x\right| \leqslant \frac{2}{n} \left[ f(2 \pi)-f(0)\right]
\]


\banswer{
	设 $n$ 为非零整数。对 $\int_{0}^{2\pi}f(x)\sin(nx)dx$ 分部积分：
	\begin{align*}
	\int_{0}^{2\pi}\!f\sin(nx)dx
	&=\left[-\frac{f(x)\cos nx}{n}\right]_{0}^{2\pi}
	+\frac1n\int_{0}^{2\pi}\!f'(x)\cos(nx)dx\\
	&=-\frac{f(2\pi)-f(0)}{n}+\frac1n\int_{0}^{2\pi}\!f'(x)\cos(nx)dx ,
	\end{align*}
	这里用到 $n$ 为整数时 $\cos 2n\pi=\cos 0=1$。

	因 $f'\geqslant0$，由牛顿--莱布尼茨公式
	\[
	\left|\int_{0}^{2\pi}\!f'(x)\cos(nx)dx\right|
	\leqslant\int_{0}^{2\pi}\!f'(x)dx=f(2\pi)-f(0) ,
	\]
	且此量非负。于是
	\[
	\left|\int_{0}^{2\pi}\!f(x)\sin(nx)dx\right|
	\leqslant\frac{2}{|n|}\left[f(2\pi)-f(0)\right].
	\]

\par
上式实际对所有非零整数 $n$ 都成立，限制 $n>0$ 即得题设中的 $2/n$ 形式。

}



\item 
(本题满分 10 分)
设 $f(x)$ 具有连续导数, 且满足 $f(x)=x+\int_{0}^{x} t f^{\prime}(x-t) d t$, 求极限 $\lim\limits _{x \rightarrow-\infty} f(x)$。


\banswer{
	\textbf{化去卷积.} 在 $\int_{0}^{x}tf'(x-t)dt$ 中作代换 $u=x-t$，得
	\[
	\int_{0}^{x}\!tf'(x-t)dt=\int_{0}^{x}\!(x-u)f'(u)du .
	\]
	于是原式化为
	\begin{align*}
	f(x)&=x+x\int_{0}^{x}\!f'(u)du-\int_{0}^{x}\!uf'(u)du\\
	&=x+x\left[f(x)-f(0)\right]-\int_{0}^{x}\!uf'(u)du .
	\end{align*}

	\textbf{定初值并求导.} 令 $x=0$ 得 $f(0)=0$。两端对 $x$ 求导：
	\[
	f'(x)=1+\left[f(x)-f(0)\right]+xf'(x)-xf'(x)=1+f(x) .
	\]

	\textbf{解方程.} 由 $f'=1+f$、$f(0)=0$ 得 $\ln|1+f|=x$，即
	\[
	f(x)=e^{x}-1 .
	\]
	（代回原式可直接验证。）故
	\[
	\lim\limits_{x\to-\infty}f(x)=\lim\limits_{x\to-\infty}\left(e^{x}-1\right)=-1 .
	\]
}


\item 
(本题满分 10 分)
设函数 $z=f(x, y)$ 满足 $\frac{\partial^{2} z}{\partial x \partial y}=x+y, f(x, 0)=x, f(0, y)=y^{2}$, 求 $f(x, y)$。

\banswer{
	\textbf{两次积分.} 把 $\frac{\partial}{\partial y}\left(\frac{\partial z}{\partial x}\right)=x+y$ 对 $y$ 积分：
	\[
	\frac{\partial z}{\partial x}=xy+\frac{y^2}{2}+\varphi(x) ;
	\]
	再对 $x$ 积分（记 $\Phi'=\varphi$）：
	\[
	z=\frac{x^2y}{2}+\frac{xy^2}{2}+\Phi(x)+\psi(y) .
	\]

	\textbf{定待定函数.} 代入两个定解条件：
	\[
	f(x,0)=\Phi(x)+\psi(0)=x,\qquad
	f(0,y)=\Phi(0)+\psi(y)=y^2 .
	\]
	两式相加得 $\Phi(x)+\psi(y)=x+y^2-\left[\Phi(0)+\psi(0)\right]$；
	再令 $x=y=0$ 知 $\Phi(0)+\psi(0)=0$，故 $\Phi(x)+\psi(y)=x+y^2$。于是
	\[
	f(x,y)=\frac{1}{2}x^{2}y+\frac{1}{2}xy^{2}+x+y^{2} .
	\]
	（验证：$\frac{\partial^2f}{\partial x\partial y}=x+y$，$f(x,0)=x$，$f(0,y)=y^2$，全部吻合。）
}



\item 
(本题满分 10 分)
求函数 $f(x, y)=(x-6)^{2}+(y+8)^{2}$ 在区域 $D=\left\{(x, y) \mid x^{2}+y^{2} \leqslant 25\right\}$ 上的最大值 和最小值。

\banswer{
	\textbf{解法一（几何意义）.} $\sqrt{f(x,y)}$ 就是点 $(x,y)$ 到定点 $A(6,-8)$ 的距离，
	而 $D$ 是以原点 $O$ 为心、$5$ 为半径的闭圆盘。由 $|OA|=\sqrt{36+64}=10>5$ 知 $A$ 在 $D$ 外，
	故 $(x,y)\in D$ 时
	\[
	|OA|-5\leqslant|AP|\leqslant|OA|+5,\quad\text{即}\quad 5\leqslant\sqrt f\leqslant15 .
	\]
	两端分别在线段 $OA$ 与圆周的交点、及其反向延长线与圆周的交点处取到，即
	\[
	f_{\min}=25\ \left(\text{于}\ (3,-4)\right),\qquad
	f_{\max}=225\ \left(\text{于}\ (-3,4)\right).
	\]

	\textbf{解法二（驻点 + 边界）.} $f_x=2(x-6)=0,\ f_y=2(y+8)=0$ 给出唯一驻点 $(6,-8)\notin D$，
	故最值必在边界 $x^2+y^2=25$ 上。令 $x=5\cos\theta,\ y=5\sin\theta$，
	\begin{align*}
	f&=25-60\cos\theta+80\sin\theta+100\\
	&=125+100\sin(\theta-\varphi),\quad \tan\varphi=\tfrac34 ,
	\end{align*}
	取值范围为 $[25,225]$，与解法一一致。
}


\newpage

\item 
(本题满分 10 分)
计算二重积分 $\iint_{D}(|x|+|y|) d x d y$, 其中积分区域 $D$ 由直线 $x=0, x+y=3, y=x-1$ 及 $y=x+1$ 所围成的区域。



\banswer{
	\textbf{定区域.} 四条直线两两相交给出四个顶点
	\[
	(0,-1),\quad(0,1),\quad(1,2),\quad(2,1),
	\]
	$D$ 是以它们为顶点的凸四边形（即夹在两条平行线 $y=x\pm1$ 之间、
	被 $x=0$ 与 $x+y=3$ 截下的部分）。在 $D$ 上恒有 $x\geqslant0$，故 $|x|=x$；
	而 $y$ 变号，需分段处理 $|y|$。按 $x$ 切片：
	$0\leqslant x\leqslant1$ 时 $x-1\leqslant y\leqslant x+1$，
	$1\leqslant x\leqslant2$ 时 $x-1\leqslant y\leqslant3-x$。

	\textbf{算 $\iint_D x\,dxdy$.}
	\begin{align*}
	\iint_{D}x\,dxdy&=\int_0^1\!\!2x\,dx+\int_1^2\!\!x(4-2x)dx\\
	&=1+\left[2x^2-\tfrac23x^3\right]_1^2=1+\tfrac43=\tfrac73 .
	\end{align*}

	\textbf{算 $\iint_D|y|\,dxdy$.} 当 $0\leqslant x\leqslant1$ 时 $x-1\leqslant0\leqslant x+1$，
	\[
	\int_{x-1}^{x+1}\!|y|dy=\frac{(1-x)^2+(1+x)^2}{2}=1+x^2 ;
	\]
	当 $1\leqslant x\leqslant2$ 时积分区间全在 $y\geqslant0$ 内，
	\[
	\int_{x-1}^{3-x}\!y\,dy=\frac{(3-x)^2-(x-1)^2}{2}=4-2x .
	\]
	故
	\[
	\iint_{D}|y|dxdy=\int_0^1\!\!\left(1+x^2\right)dx+\int_1^2\!\!(4-2x)dx=\tfrac43+1=\tfrac73 .
	\]

	\textbf{合并.} $I=\dfrac73+\dfrac73=\dfrac{14}{3}\approx4.667$（数值积分给出 $4.6667$，吻合）。
}


\item 
(本题满分 10 分)
计算曲线积分 $I=\oint_{L} \frac{x d y-y d x}{4 x^{2}+y^{2}}$, 其中 $L$ 是以 $(1,0)$ 为中心, 半径为 $R$ 的圆周 $(R>0, R \neq 1)$, 取逆时针方向。

\banswer{
	记 $P=\dfrac{-y}{4x^2+y^2}$，$Q=\dfrac{x}{4x^2+y^2}$。在除原点外的区域内
	\[
	\frac{\partial Q}{\partial x}=\frac{y^2-4x^2}{\left(4x^2+y^2\right)^2}
	=\frac{\partial P}{\partial y} ,
	\]
	故被积表达式在 $\mathbb{R}^2\setminus\{O\}$ 内是闭形式，唯一的奇点是原点。
	$L$ 的圆心为 $(1,0)$，$|O-(1,0)|=1$，于是分两种情形。

	\textbf{情形一：$0<R<1$.} 原点在 $L$ 之外，$L$ 所围闭区域内 $P,Q$ 有连续偏导，
	由格林公式 $I=\iint\left(Q_x-P_y\right)dxdy=0$。

	\textbf{情形二：$R>1$.} 原点在 $L$ 之内。取充分小的 $\varepsilon>0$，
	以椭圆 $\ell_\varepsilon:\ 4x^2+y^2=\varepsilon^2$（逆时针）挖去原点；
	在 $L$ 与 $\ell_\varepsilon$ 之间的复连通区域上用格林公式得 $I=\oint_{\ell_\varepsilon}$。
	令 $x=\frac{\varepsilon}{2}\cos t,\ y=\varepsilon\sin t$，则
	\[
	x\,dy-y\,dx=\frac{\varepsilon^2}{2}dt,\qquad 4x^2+y^2=\varepsilon^2 ,
	\]
	故
	\[
	I=\int_0^{2\pi}\frac{1}{2}dt=\pi .
	\]

	\textbf{结论.} $0<R<1$ 时 $I=0$；$R>1$ 时 $I=\pi$（与 $R$ 无关）。
	数值积分取 $R=0.5,0.9$ 得 $0$，取 $R=1.5,3,10$ 得 $3.14159\cdots$，与结论一致。
}



\item 
(本题满分 10 分)
设函数 $f(x)$ 在 $(-\infty,+\infty)$ 上有定义，在 $x=0$ 的某个邻域内有一阶连续的导数, 且 $\lim\limits _{x \rightarrow 0} \frac{f(x)}{x}=a>0$, 证明:
$\sum\limits_{n=1}^{\infty}(-1)^{n} f\left(\frac{1}{n}\right)$ 收敛，而 $\sum\limits_{n=1}^{\infty} f\left(\frac{1}{n}\right)$ 发散。



\banswer{
	\textbf{预备.} 由 $\lim\limits_{x\to0}\frac{f(x)}{x}=a$ 有限，且 $f$ 在 $x=0$ 附近连续，
	必有 $f(0)=0$，从而
	\[
	f'(0)=\lim\limits_{x\to0}\frac{f(x)-f(0)}{x}=a>0 .
	\]
	又 $f'$ 在 $0$ 的某邻域内连续，由保号性存在 $\delta>0$ 使 $f'(x)>0$（$|x|<\delta$），
	于是 $f$ 在 $(0,\delta)$ 上严格递增。取 $N>\frac1\delta$，则当 $n\geqslant N$ 时
	$0<\frac1{n+1}<\frac1n<\delta$，故
	\[
	f\left(\tfrac1n\right)>f\left(\tfrac1{n+1}\right)>f(0)=0 ,
	\]
	即 $\left\{f\left(\frac1n\right)\right\}_{n\geqslant N}$ 单调递减且为正；
	又由连续性 $f\left(\frac1n\right)\to f(0)=0$。

	\textbf{交错级数收敛.} 由上，从第 $N$ 项起 $\sum(-1)^nf\left(\frac1n\right)$ 满足
	莱布尼茨判别法的全部条件，故收敛；前面有限项不影响敛散性，
	所以 $\sum\limits_{n=1}^{\infty}(-1)^{n}f\left(\frac1n\right)$ 收敛。

	\textbf{正项级数发散.} 由
	$\lim\limits_{n\to\infty}\dfrac{f\left(\frac1n\right)}{\frac1n}=a>0$
	及比较判别法的极限形式，$\sum f\left(\frac1n\right)$ 与调和级数 $\sum\frac1n$ 同敛散；
	而 $\sum\frac1n$ 发散，故 $\sum\limits_{n=1}^{\infty}f\left(\frac1n\right)$ 发散。

	\par\vspace{1mm}\noindent
	\textbf{注：}{\kaishu
	交错级数一段不能只用 $f\left(\frac1n\right)\sim\frac an$ 加“$\sum\frac{(-1)^n}{n}$ 收敛”搪塞：
	等价无穷小只给出 $f\left(\frac1n\right)=\frac an+o\left(\frac1n\right)$，
	而 $o\left(\frac1n\right)$ 的交错和未必收敛。必须像上面那样用单调性（即 $f'>0$）走莱布尼茨。
	}\par
}


\item 
(本题满分 10 分)
将函数 $f(x)=x^{2}(-\pi \leqslant x \leqslant \pi)$ 展开成以 $2 \pi$ 为周期的傅里叶级数。


\banswer{
	$f(x)=x^2$ 在 $[-\pi,\pi]$ 上是偶函数，故 $b_n=0$，只需算余弦系数。

	\textbf{常数项.}
	\[
	a_0=\frac1\pi\int_{-\pi}^{\pi}\!x^2dx=\frac{2}{\pi}\cdot\frac{\pi^3}{3}=\frac{2\pi^2}{3},
	\qquad \frac{a_0}{2}=\frac{\pi^2}{3} .
	\]

	\textbf{余弦系数.} 两次分部积分（$n\geqslant1$）：
	\begin{align*}
	a_n&=\frac2\pi\int_0^{\pi}\!x^2\cos nx\,dx
	=\frac2\pi\left(\left[\frac{x^2\sin nx}{n}\right]_0^{\pi}-\frac2n\int_0^{\pi}\!x\sin nx\,dx\right)\\
	&=-\frac{4}{n\pi}\int_0^{\pi}\!x\sin nx\,dx
	=-\frac{4}{n\pi}\cdot\left(-\frac{\pi\cos n\pi}{n}\right)
	=\frac{4(-1)^{n}}{n^{2}} .
	\end{align*}

	\textbf{结论.} $f$ 以 $2\pi$ 为周期延拓后处处连续（端点 $f(-\pi)=f(\pi)=\pi^2$），
	由狄利克雷定理级数处处收敛于 $f$：
	\[
	x^{2}=\frac{\pi^{2}}{3}+\sum\limits_{n=1}^{\infty}\frac{4(-1)^{n}}{n^{2}}\cos nx,
	\qquad -\pi\leqslant x\leqslant\pi .
	\]

	\par\vspace{1mm}\noindent
	\textbf{注：}{\kaishu
	取 $x=\pi$ 得 $\sum\limits_{n=1}^{\infty}\frac1{n^2}=\frac{\pi^2}{6}$，
	取 $x=0$ 得 $\sum\limits_{n=1}^{\infty}\frac{(-1)^{n-1}}{n^2}=\frac{\pi^2}{12}$，
	可作为结果的校验。
	}\par
}


\item 
(本题满分 10 分)
已知函数 $f(x)$ 具有连续的二阶导数, $f(0)=0, f^{\prime}(0)=\frac{1}{3}$, 且对任意光滑有向封闭 曲面 $\Sigma$都有 $\oint_{\Sigma} e^{x}\left(f^{\prime}(x) d y d z-2 y f(x) d z d x-z d x d y\right)=0$, 求函数 $f(x)$ 的表达式。


\banswer{
	\textbf{由高斯公式导出方程.} 记
	\[
	\bm{F}=\left(e^{x}f'(x),\ -2ye^{x}f(x),\ -ze^{x}\right).
	\]
	对任意光滑有向闭曲面 $\Sigma$（设其所围区域为 $\Omega$），高斯公式给出
	$\oiint_{\Sigma}\bm{F}\cdot d\bm{S}=\pm\iiint_{\Omega}\operatorname{div}\bm{F}\,dV$。
	要它对一切 $\Sigma$ 都为零，必须 $\operatorname{div}\bm{F}\equiv0$。而
	\begin{align*}
	\operatorname{div}\bm{F}&=e^{x}\left[f'(x)+f''(x)\right]-2e^{x}f(x)-e^{x}\\
	&=e^{x}\left[f''+f'-2f-1\right] ,
	\end{align*}
	故 $f$ 满足
	\[
	f''+f'-2f=1 .
	\]

	\textbf{解方程.} 特征方程 $r^2+r-2=0$ 的根为 $r=1,-2$；
	常数项 $1$ 对应特解 $f^{*}=-\frac12$。故
	\[
	f(x)=C_1e^{x}+C_2e^{-2x}-\frac12 .
	\]

	\textbf{定常数.} 由 $f(0)=0$ 得 $C_1+C_2=\frac12$；
	由 $f'(0)=\frac13$ 得 $C_1-2C_2=\frac13$。解得 $C_1=\frac49,\ C_2=\frac1{18}$，即
	\[
	f(x)=\frac{4}{9}e^{x}+\frac{1}{18}e^{-2x}-\frac{1}{2} .
	\]
	（验证：$f''+f'-2f=1$，$f(0)=0$，$f'(0)=\frac13$，全部满足。）
}



\item 
(本题满分 10 分)
已知平面过点 $ (1,2,3) $，它在 $ x $ 正半轴和 $ y $ 轴正半轴上的截距相等，问当该平面在
三个坐标轴上的截距分别为何值时，它与三个坐标轴所围成的立体体积取极值，并写出此时的平面方程。


%群友反馈：原题目有问题，最小值为零，两种改法：①将$ z $轴的截距限制为正；②将最值改为极值。


\banswer{
	\textbf{建模.} 设 $x$、$y$ 轴上的截距同为 $a>0$，$z$ 轴上的截距为 $c$，则平面为
	$\frac xa+\frac ya+\frac zc=1$。过 $(1,2,3)$ 给出约束
	\[
	\frac1a+\frac2a+\frac3c=1,\qquad\text{即}\qquad \frac3a+\frac3c=1 .
	\]
	立体体积 $V=\frac16a\cdot a\cdot c=\frac{a^2c}{6}$。

	\textbf{消元求极值.} 由约束解出 $c=\frac{3a}{a-3}$（取 $a>3$ 以保证 $c>0$），代入得
	\[
	V(a)=\frac{a^2}{6}\cdot\frac{3a}{a-3}=\frac{a^3}{2(a-3)} ,
	\]
	\[
	V'(a)=\frac{a^{2}(2a-9)}{2(a-3)^{2}} .
	\]
	在 $a>3$ 上 $V'<0$（$3<a<\frac92$）、$V'>0$（$a>\frac92$），故 $a=\frac92$ 处取极小值。此时
	\[
	c=\frac{3\cdot\frac92}{\frac92-3}=9,\qquad
	V_{\min}=\frac{\left(\frac92\right)^2\cdot9}{6}=\frac{243}{8} .
	\]

	\textbf{结论.} $x$ 轴、$y$ 轴、$z$ 轴上的截距分别为 $\frac92$、$\frac92$、$9$ 时体积取极（小）值
	$\frac{243}{8}$，此时平面方程为
	\[
	\frac{x}{9/2}+\frac{y}{9/2}+\frac{z}{9}=1,\qquad\text{即}\qquad 2x+2y+z=9 .
	\]

	\par\vspace{1mm}\noindent
	\textbf{注：}{\kaishu
	原答案把 $x$、$y$ 截距写成 $\frac29$，与它自己给出的平面方程 $\frac{x}{9/2}+\frac{y}{9/2}+\frac z9=1$
	自相矛盾，应为 $\frac92$；数值搜索也给出 $a=\frac92$，$V=\frac{243}{8}=30.375$。

	另外题面说“取极值”而非“取最值”是必要的：若允许 $0<a<3$，则 $c=\frac{3a}{a-3}<0$，
	体积 $|V|=\frac{a^3}{2(3-a)}$ 当 $a\to0^{+}$ 时趋于 $0$ 而取不到，故全局最小值不存在。
	}\par
}


\item 
(本题满分 10 分)
设函数 $f(x)$ 在 $[0, c]$ 上可导, $f^{\prime}$ 单调递减且 $f(0)=0$, 求证: 对于 $0 \leqslant a \leqslant b \leqslant a+b \leqslant c$ ，都有 $f(a+b) \leqslant f(a)+f(b)$。

\banswer{
	若 $a=0$，由 $f(0)=0$ 知 $f(a+b)=f(b)=f(a)+f(b)$，结论取等号成立。以下设 $a>0$。

	\textbf{两次拉格朗日中值定理.} 在 $[0,a]$ 上对 $f$ 用中值定理：存在 $\xi\in(0,a)$ 使
	\[
	f(a)=f(a)-f(0)=f'(\xi)\,a .
	\]
	在 $[b,a+b]$ 上（长度同为 $a$，且 $a+b\leqslant c$）再用一次：存在 $\eta\in(b,a+b)$ 使
	\[
	f(a+b)-f(b)=f'(\eta)\,a .
	\]

	\textbf{比较.} 由 $0<\xi<a\leqslant b<\eta$ 及 $f'$ 单调递减得 $f'(\xi)\geqslant f'(\eta)$，于是
	\[
	f(a)=f'(\xi)a\geqslant f'(\eta)a=f(a+b)-f(b) ,
	\]
	即
	\[
	f(a+b)\leqslant f(a)+f(b) .
	\]

	\par\vspace{1mm}\noindent
	\textbf{注：}{\kaishu
	条件 $a\leqslant b$ 正是为了保证两个中值点满足 $\xi<\eta$；
	$f'$ 递减（即 $f$ 为凹函数）加上 $f(0)=0$，本质上说明 $\frac{f(x)}{x}$ 递减，
	这正是次可加性 $f(a+b)\leqslant f(a)+f(b)$ 的来源。
	}\par
}




\end{enumerate}


